\subsection{3. Exakter Vierzustandskanal mit dem vollständigen nativen
W}\label{exakter-vierzustandskanal-mit-dem-vollstuxe4ndigen-nativen-w}

Die gepinnte Quelle ist

\codeword{universalraum-native-ground-response-20260915/ground\_replay/outputs/simple\_core/spinor\_tensors.npz},

SHA-256
\codeword{3f00a0892157a691e4ef26920c3702772c7bf76a250fad98ff04a079bb64b763}.

Die zwei relevanten Quellenzeilen sind

\[\adjustbox{max width=\linewidth}{$\displaystyle 
\begin{aligned}
P_0={}&-f_{57}f_4+f_{56}f_5+f_{53}f_8-f_{52}f_9
       -f_{45}f_{16}+f_{44}f_{17}+f_{33}f_{28}-f_{32}f_{29},\\
P_1={}&-f_{58}f_4+f_{56}f_6+f_{54}f_8-f_{52}f_{10}
       -f_{46}f_{16}+f_{44}f_{18}+f_{34}f_{28}-f_{32}f_{30}.
\end{aligned}
$}\]

Die Konvention ist \(P_A=\sum_{i<j}W_{A,ij}f_jf_i\). Das gemeinsame
aktive Fermion ist s=4; die Endpunktmarken sind 57 und 58.

Präpariert werden sieben unveränderte Pauli-Blocker

\[\adjustbox{max width=\linewidth}{$\displaystyle 
S=\{8,16,28,32,44,52,56\}.
$}\]

Sei \(|S\rangle\) der geordnet erzeugte Fockzustand dieser Moden. Die
vier Basiszustände sind

\[\adjustbox{max width=\linewidth}{$\displaystyle 
\begin{array}{c|c|c}
\text{Zustand}&\text{besetzte Fermionmoden}&\text{Bosonen}\\\hline
L&S\cup\{4,57\}&0\\
B_0&S&1\text{ in Kanal }0\\
B_1&S&1\text{ in Kanal }1\\
R&S\cup\{4,58\}&0
\end{array}
$}\]

Jeder Zustand besitzt die native Gesamtladung \(N=N_f+2N_b=9\).

Zugelassen wird nun genau eine zusätzliche Hamiltonoperation,

\[\adjustbox{max width=\linewidth}{$\displaystyle 
H_J=J(b_1^\dagger b_0+b_0^\dagger b_1).
$}\]

Alle ursprünglichen 60 nativen Paaroperatoren bleiben erhalten:

\[\adjustbox{max width=\linewidth}{$\displaystyle 
H=\Delta N_b+g\sum_{A=0}^{59}(b_A^\dagger P_A+P_A^\dagger b_A)+H_J.
$}\]

Der Prüfer wendet sämtliche 480 signierten Paarkanäle einschließlich
aller CAR-Vorzeichen und der Bosonfaktoren auf jeden der vier Zustände
an. Es entstehen \textbf{keine Zustände außerhalb ihres linearen
Spanns}. Die Pauli-Blocker verhindern insbesondere die sieben
unerwünschten Rückpaarungen jedes aktiven Bosonkanals.

Es ergibt sich exakt

\[\adjustbox{max width=\linewidth}{$\displaystyle 
\boxed{
H_4=\begin{pmatrix}
0&g&0&0\\
g&\Delta&J&0\\
0&J&\Delta&g\\
0&0&g&0
\end{pmatrix}_{(L,B_0,B_1,R)}.
}
$}\]

Ohne Mischer, J=0, zerfällt die Matrix in zwei getrennte 2×2-Blöcke.
Dann ist \(\langle R|e^{-itH}|L\rangle\) zu jeder Zeit exakt null. Mit J
ungleich null beginnt die Transferamplitude bei dritter Ordnung; die
Wahrscheinlichkeit hat den führenden Term \(g^4J^2t^6/36\).

\textbf{Die Quelle enthält hier also die beiden Endstücke und eine exakt
funktionierende Pauli-Blockierung, aber nicht die geprüfte
Bosonverbindung als bereits verfügbare Operation.}

\subsection{4. Strenge vollständige Transfergrenze im tatsächlichen
W-Zeugen}\label{strenge-vollstuxe4ndige-transfergrenze-im-tatsuxe4chlichen-w-zeugen}

Unter Spiegelung L↔R und B₀↔B₁ zerfällt H₄ in

\[\adjustbox{max width=\linewidth}{$\displaystyle 
H_+=\begin{pmatrix}0&g\\g&\Delta+J\end{pmatrix},\qquad
H_-=\begin{pmatrix}0&g\\g&\Delta-J\end{pmatrix}.
$}\]

Wähle den vorhandenen Prüfquotienten \(g/\Delta=1/20\) und die
ausdrücklich neu gesetzte Mischerstärke

\[\adjustbox{max width=\linewidth}{$\displaystyle 
J=\Delta-\frac{2g}{\sqrt3}
 =\Delta\left(1-\frac1{10\sqrt3}\right),
\qquad t=\frac{\pi\sqrt3}{g}=\frac{20\pi\sqrt3}{\Delta}.
$}\]

Da \(0<J<\Delta\), bleiben die Eigenfrequenzen des Bosonmischers
\(\Delta\pm J\) positiv. Es wurde \textbf{keine Nullfrequenz durch
exakte Aufhebung von \(\Delta N_b\)} erzeugt. Auch auf dem ganzen
Fockraum bleibt die hinzugefügte Hamiltonfamilie nach unten
beschränkbar: Das Bosonquadrat ist strikt positiv und die endlichen
Fermionoperatoren können durch quadratische Ergänzung kontrolliert
werden. Das bedeutet nicht, dass ihr Grundzustand noch der unveränderte
native Grundzustand ist.

Für H₋ ist \(\Delta-J=2g/\sqrt3\). Zu der gewählten Zeit kehrt dessen
Endpunktkomponente exakt mit Phase −1 zurück, ohne mittlere
Restbesetzung.

Für H₊ setze

\[\adjustbox{max width=\linewidth}{$\displaystyle 
a=\Delta-\frac g{\sqrt3},\quad
\omega=\sqrt{a^2+g^2},\quad
w_- =\frac{1-a/\omega}{2}.
$}\]

Die Endpunktamplitude lautet

\[\adjustbox{max width=\linewidth}{$\displaystyle 
A_+=(1-w_-)e^{i(\omega-a)t}+w_-e^{-i(\omega+a)t}.
$}\]

Mit \(\omega-a\le g^2/(2a)\), \(w_-\le g^2/(4a^2)\),
\(\cos u\ge1-u^2/2\) und \(|A_+|^2\ge1-g^2/a^2\) folgt

\[\adjustbox{max width=\linewidth}{$\displaystyle 
\begin{aligned}
p_{L\to R}
 &=\left|\frac{A_++1}{2}\right|^2\\
 &\ge1-\frac{g^4t^2}{16a^2}-\frac{g^2}{2a^2}\\
 &=1-\frac{3\pi^2+8}{6400(a/\Delta)^2}.
\end{aligned}
$}\]

Aus \(\sqrt3>17/10\) folgt \(a/\Delta>33/34\). Somit ergibt sich rein
rational

\[\adjustbox{max width=\linewidth}{$\displaystyle 
\boxed{
p_{L\to R}>
1-\frac{3(355/113)^2+8}{6400(33/34)^2}
=\frac{2009992727}{2022609600}
>0.9937620819>0.993.
}
$}\]

Die Aussage ist eine \textbf{analytische Schranke der vollständigen
Dynamik auf einem exakt invarianten Quellen-Unterraum}. Es gibt hier
keinen numerisch abgeschnittenen Restzustandsraum. Die
64-Fermion/60-Boson-Quelle ist nicht durch vier frei erfundene
Matrixeinträge ersetzt worden: Ihre exakte Einschränkung wurde zuerst
vollständig geprüft.

Die Zeit in Einheiten \(\hbar=1\) ist ungefähr \(108.83/\Delta\). Diese
numerische Orientierung und die gesetzte Mischerstärke sind keine
vorhergesagten Naturkonstanten.

\subsection{5. Was dies für die fundamentale Suche
ändert}\label{was-dies-fuxfcr-die-fundamentale-suche-uxe4ndert}

Der konstruktive Anschluss ist klein: \textbf{Paarumwandlung → geteilte
Zwischenressource → Paar-Rückumwandlung} kann einen
Einfermion-Endpunktwechsel bewirken. Für die tatsächliche W-Quelle
braucht man dazu weder einen neuen direkten Fermion-Hopping-Term noch
die Kontrolle jedes einzelnen der 480 Paarmonome. Ein einzelner
Bosonmischer plus eine spezielle Pauli-blockierte Präparation reicht im
belegten internen Zeugen.

Offen bleiben aber genau die folgenden Herkunftsfragen:

\begin{enumerate}
\def\labelenumi{\arabic{enumi}.}
\tightlist
\item
  \textbf{Operationssatz:} Ist der reine Mischer
  \(b_1^\dagger b_0+\mathrm{h.c.}\) tatsächlich verfügbar? Eine
  simultane Fermion-und-Boson-Symmetrie ist nicht automatisch eine
  unabhängige Bosonoperation.
\item
  \textbf{Präparation:} Wie wird der N=9-Blockerzustand mit dem
  zugelassenen Operationssatz erzeugt? Ladungserhaltende native
  Evolution präpariert ihn nicht aus dem N=64-Grundzustand.
\item
  \textbf{Operationaler Raum:} Sind die Endpunktmarken 57 und 58
  überhaupt verschiedene räumliche Teile oder nur interne Marken
  derselben Bank? Die Rechnung allein liefert keine Raumposition.
\item
  \textbf{Gemeinsamer Grundzustand:} Besteht ein entsprechender
  Mechanismus für die ursprüngliche Entnahmeantwort des nativen
  N=64-Grundzustands, statt für diesen besonders präparierten Zeugen?
\item
  \textbf{Zwei Banken:} Eine Bosonverbindung zwischen wirklich
  unabhängigen Banken erhält deren lokale Fermionparitäten. Dieser
  interne Zeuge hebt jene Schranke nicht auf. Dafür wäre eine
  ursprünglich geteilte Fermionstruktur oder eine andere insgesamt
  gerade, lokal ungerade Quelloperation zu konstruieren.
\end{enumerate}

T1-T8, das relativistische Feldwörterbuch und die Herkunft der Raumzeit
sind damit nicht geschlossen. Das Ergebnis verkleinert eine konkrete
Suche: Statt eines völlig beliebigen neuen Fermionlinks kann jetzt ein
bestimmter Bosonmischer mitsamt seinem Quellen- und Präparationsvertrag
geprüft werden. Die nachfolgende Quellenprüfung weist diesen Mischer
jedoch ausdrücklich \textbf{nicht} als Synthese der bisher gewährten
Kontrollen aus.

\subsection{6. Anschluss des neuen Chart-/Glue-Vorschlags und exakter
Quellen-No-go}\label{anschluss-des-neuen-chart-glue-vorschlags-und-exakter-quellen-no-go}

Der neue Nutzeranhang vom 15. September 2026 wurde vollständig gelesen:

\codeword{/Users/stefanhamann/.codex/attachments/59dc0059-8914-48ca-953d-85933f66e00b/pasted-text.txt},

SHA-256
\codeword{1a75e84b28868dd682e4b2337d6546f275f49321e489b3b23af4754a844d1ac4}.

Die Hypothese, Banken zunächst als überlappende lokale Beschreibungen
einer gemeinsamen CAR-Struktur zu untersuchen, passt zum hier
konstruierten gemeinsamen Fermion-/Boson-Zwischenraum. \textbf{Sie löst
den Übergang von Überlappung zu Dynamik aber nicht automatisch.} Bereits
zwei nichtorthogonale lokale Moden \(f_A=f_1\),
\(f_B=\cos\theta f_1+\sin\theta f_2\) bei H=0 besitzen ein
nichtverschwindendes Kreuz-Antikommutator
\(\{f_A,f_B^\dagger\}=\cos\theta\), obwohl überhaupt keine
Zustandsentwicklung stattfindet. Eine Änderung der Beschreibung und ein
dynamischer Transfer müssen getrennt geprüft werden.

\subsubsection{6.1 Eine nötige Korrektur des vorgeschlagenen
Paritätstests}\label{eine-nuxf6tige-korrektur-des-vorgeschlagenen-parituxe4tstests}

Der Anhang fordert für überlappende Charts einen Intertwiner, der lokal
Paritäten ändert, aber mit \(\Pi_A\Pi_B\) kommutiert. \textbf{Bei
Überlappung ist dieses Produkt nicht automatisch die Gesamtparität.} Im
kleinsten gemeinsamen-Moden-Zeugen

\[\adjustbox{max width=\linewidth}{$\displaystyle 
A=\{x,s\},\qquad B=\{s,y\}
$}\]

gilt

\[\adjustbox{max width=\linewidth}{$\displaystyle 
\Pi_A\Pi_B=(-1)^{n_x+n_y},
\qquad
\Pi_{\rm global}=(-1)^{n_x+n_s+n_y}.
$}\]

Die gemeinsame Mode s wurde im Produkt zweimal gezählt und fällt heraus.
Die Paarumwandlung \(b^\dagger s f_x\) ist bezüglich
\(\Pi_{\rm global}\) gerade, kommutiert aber nicht mit \(\Pi_A\Pi_B\).
Der vorgeschlagene Kill-Test würde hier einen legitimen global geraden
Überlappungsmechanismus fälschlich aussortieren. Alle vier Identitäten
wurden auf dem vollständigen Drei-Fermion-CAR-Raum exakt geprüft.

\textbf{Der korrigierte Test muss die globale Parität aus der
gemeinsamen CAR-Darstellung selbst verwenden}, nicht aus einem
ungeprüften Produkt lokaler Paritäten. Zusätzlich benötigt er eine
definierte Anfangspräparation, dynamisch unterschiedliche
Endpunktbeobachtungen und einen aus der Quelle stammenden Generator.
Passive Chartwechsel allein reichen nicht.

\subsubsection{6.2 Tatsächlicher endlicher Clock statt unterstellter
voller
Gruppe}\label{tatsuxe4chlicher-endlicher-clock-statt-unterstellter-voller-gruppe}

Die Prüfung lädt den gepinnten ursprünglichen Clock-Konstruktor über

\codeword{sources/repo/experiments/theory-contracts/universalraum-native-operations-ground-response-20260915/common.py},

SHA-256
\codeword{2cc97522457ecc7e6774d96a25b2581b5d67cd051ae35e28a26f0e98e2adb994}.

Der Konstruktor kontrolliert seine ursprünglichen Quellenpins und die
exakte W-Kovarianz. Er liefert tatsächlich

\[\adjustbox{max width=\linewidth}{$\displaystyle 
p=(2,0,1,4,3),\qquad s_{\rm Clock}=-1,
$}\]

und auf Bosonmoden \(A=6k+c\)

\[\adjustbox{max width=\linewidth}{$\displaystyle 
G_B e_{6k+c}=-e_{6(p(k\bmod5)+5\lfloor k/5\rfloor)+c}.
$}\]

Es wird \textbf{nicht} vorausgesetzt, dass der Clock der volle
Spin(10)×SU(4)-Operationssatz ist. Sein endlicher, direkt reproduzierter
Lift reicht für den folgenden Ausschluss.

Schreibe \(M_{ij}=|i\rangle\langle j|+|j\rangle\langle i|\) auf dem
Boson-Einteilchenraum. Exakt gilt

\[\adjustbox{max width=\linewidth}{$\displaystyle 
G_BM_{01}G_B^\dagger=M_{12,13},\quad
G_BM_{12,13}G_B^\dagger=M_{6,7},\quad
G_BM_{6,7}G_B^\dagger=M_{01}.
$}\]

Insbesondere

\[\adjustbox{max width=\linewidth}{$\displaystyle 
\|[G_B,M_{01}]\|_F^2=4\ne0.
$}\]

Die native W-Kovarianz impliziert \([G,H]=[G,N_b]=0\) für den
gemeinsamen Fock-Lift G. Auch ein gewährter getrennter Gesamt-Casimir
von Spin(10) oder SU(4) kommutiert mit diesem tatsächlichen G. Daher
liegt

\[\adjustbox{max width=\linewidth}{$\displaystyle 
\operatorname{Alg}(H,N_b,G,C_{\rm Spin},C_{\rm Color})\subset\{G\}'
$}\]

und der spezifische Einzelmischer \(b_1^\dagger b_0+\mathrm{h.c.}\)
liegt \textbf{nicht} in dieser erzeugten Algebra. Auf den endlichen
Ladungssektoren gibt es dabei keine Domänenprobleme; für die volle
Hilbert-Darstellung formuliert man dieselbe Aussage mit den erzeugten
beschränkten Zeitentwicklungen und Spektraloperationen.

\subsubsection{6.3 Clock-Mittelung ist möglich, aber der Farb-Cartan
sperrt sie
weiterhin}\label{clock-mittelung-ist-muxf6glich-aber-der-farb-cartan-sperrt-sie-weiterhin}

Der Clock-Ausschluss allein wäre zu schwach. Die Orbitsumme

\[\adjustbox{max width=\linewidth}{$\displaystyle 
M_{\rm orb}=M_{01}+M_{12,13}+M_{6,7}
$}\]

kommutiert exakt mit \(G_B\). Ihre beiden zusätzlichen Linkblöcke wirken
auf den vier Zuständen aus Abschnitt 3 null. Ein neu zugelassener
globaler Bosonmischer \(J b^\dagger M_{\rm orb}b\) erzeugt deshalb
\textbf{denselben} exakt invarianten Vierzustandskanal und dieselbe
Transfergrenze. Ein neuer globaler Zusammenhang muss also nicht zwingend
den endlichen Clock brechen.

Es gibt jedoch einen zweiten, stärkeren gemeinsamen Erhaltungssatz. Aus
dem tatsächlichen Quell-Gewichtswörterbuch wählen wir die zweite
SU(4)-Cartankomponente und definieren den gemeinsamen Ladungsoperator

\[\adjustbox{max width=\linewidth}{$\displaystyle 
Q=\sum_r q_r n_{f,r}+\sum_A q_A n_{b,A}.
$}\]

Jeder ursprüngliche Paarterm erfüllt exakt \(q_i+q_j=q_A\); dies wird
für alle 480 Quellenpaare geprüft. Der tatsächliche Clock wirkt auf dem
Spinindex und lässt die SU(4)-Farbkomponente unverändert. Daher

\[\adjustbox{max width=\linewidth}{$\displaystyle 
[Q,H]=[Q,N_b]=[Q,G]=0.
$}\]

Die getrennten Gesamt-Casimire kommutieren ebenfalls mit diesem Cartan.
Diese Aussage erfordert \textbf{keine} Verfügbarkeit aktiver
kontinuierlicher SU(4)-Kontrollen.

Auf den beiden Bosonkanälen gilt dagegen

\[\adjustbox{max width=\linewidth}{$\displaystyle 
q_0=0,\qquad q_1=2,
$}\]

und auf dem ganzen Boson-Einteilchenraum

\[\adjustbox{max width=\linewidth}{$\displaystyle 
\|[Q_B,M_{01}]\|_F^2=8,\qquad
\|[Q_B,M_{\rm orb}]\|_F^2=24.
$}\]

\textbf{Damit sind sowohl der Einzelmischer als auch seine
Clock-invariante Orbitsumme aus dem genannten Operationssatz
ausgeschlossen.} Die Mittelung beseitigt das Clock-Hindernis, nicht den
separaten Farb-Cartan-Erhaltungssatz.

Die tatsächlichen Cartanladungen im Vierzustandsraum sind

\[\adjustbox{max width=\linewidth}{$\displaystyle 
(Q_L,Q_{B_0},Q_{B_1},Q_R)=(7,7,9,9).
$}\]

Dieser Zeuge verändert also eine innere SU(4)-Quantenzahl. Er ist noch
deutlicher als ein bloßer Markenwechsel vom Transport \textbf{derselben}
niedrigenergetischen Fermionmode zwischen zwei räumlichen Banken zu
unterscheiden.

